% Nouveau EGA V (cote 25): archive cote 25, site page 155 (PDF page 156).
% typescript. Draft compared: none.
% AI-transcribed by Claude Opus 5.5 (claude-opus-5-5, adaptive thinking, effort high), Message Batches API, batch ega25b (msgbatch_01A7n5YuXubUPNK1zsepi8RG), 2026-11-13; not yet checked by eye.
\scanpage{25/155}
\mspage{16}

inversibles). Posant $V^2 = V(D)$ , qui est un sous-préschéma fermé de
$V^1-V''$ , on trouve donc que l'ensemble des \xout{\unsure{points}} zéros singuliers
ordinaires rel.\xout{Y}S n'est autre que $V^1 - V'' - V^2$ , i.e. l'ensemble des
zéros supersinguliers \xout{\unsure{sur}} rel à S est $V''\ \ V^2$ .

\margin{$\|$} Au moment de l'introduction de V", en affirmant que $\underline{W}$ et $\underline{K}$ sont
localement libres sur $V^1-V''$, j'ai utilisé en f'ait général, que je dégage
ici en lemme (où $V^1-V''$ est devenu X):

\uline{Lemme} Soient $v\colon \underline{E} \to \underline{F}$ un homomorphisme de Modules localement libres
sur le préschéma X, \xout{\unsure{et soit}} m un entier $>0$, supposons que $\bigwedge^{m} u$ soit
nul. \xout{\unsure{(ce qui implique que le rang de u en chaque point est $\leq$ m) et que}} \inserted[typed]{Alors}
$\bigwedge^{m-1} u$ ne s'annule en aucun point (du point de vue fibres réduites, bien
sur) \xout{\unsure{xxxAlorsx}} \inserted[typed]{sss} u($\underline{E}$) est un sous-Module localement facteur direct de $\underline{F}$,
localement libre de rang m-1, \hand{ou encore ssi u est de rang \unsure{m}-1 en}
\hand{chaque pt de X.}

Il y a lieu de résumer ces constructions dans une proposition
qui généralise prop.6 . Il faut voir comment on peut généraliser de même
prop.7. Pour ceci, on suppose que Y lui-même est lisse sur un préschéma
S , \inserted[typed]{de dimension relative N,} donc X est lisse sur S, de dimension relative N+d\unsure{\hand{$^*$}}; où d est celle de X
sur Y. On supposera de plus avec les notations précédentes $V''=\emptyset$, i.e.
$\underline{W}$ localement libre de rang d-m+1 \inserted[typed]{(c'est tjrs le cas si m=1 !)}. Enfin, on supposera les sous-présché-%
mas fermés \xout{\unsure{Vx}} \inserted[typed]{$V^o$et $V^1$} de X lisses \uline{sur} S , de dimension relative \xout{minima} en chaque
point minima (pour les données d,N,m) , \xout{\unsure{exxqmx}} c'est à dire, comme
on constate aussitôt, de dimensions relatives (N+d) - m et
(N+d) - m - (d-m +1) = N - 1 . (Du point de vue rédaction, le mieux
semble de poser brutalement ces nombres, peut être plutôt en termes
de codimension fibre par fibre, celle de $V^o$ dans X étant m, celle de
$V^1$ dans $V^o$ étant d-m+1 ; et de signaler en remarque la justification
de ce choix par le principe \unsure{du|de} la dimension minima ou dimension \xout{générique}

\margin{155}
